\section{Descent to the original parameter}
\label{sec:original-parameter}

The scaled coordinate $t=d^dc^{d-1}$ is the natural coordinate on the
unicritical moduli quotient.  Returning to the original $c$-line
introduces a second, genuine Kummer descent.  We record the resulting
factorization, both to delimit the role of the scaled coordinate and to
show that this descent can still be made exact.

Retain
\[
 e=d-1,\qquad \gamma=(m,e),\qquad
 q=e/\gamma,\qquad a=m/\gamma,
\]
so $(a,q)=1$.  Choose $s\in\overline\Q$ with $s^q=d$ and put
\[
 B_0=\Q(s),\qquad \rho=s^aP.
 \numberthis\label{eq:rho-definition}
\]
Let $\mathcal C_{d,m}/\Q$ be the normalization of the projective
multiplier curve $M_{d,m}(c,x)=0$; thus
$\Q(\mathcal C_{d,m})=\Q(c,x)$.  By
Lemma~\ref{lem:cycle-scaling-quotients} and Galois descent, the cycle
product $P$ belongs to $\Q(c,x)$ and satisfies $x=d^mP^e$.  Since
$ae=mq$, one has
\[
 \rho^e=s^{ae}P^e=d^mP^e=x.
\]
Moreover,
\[
 B_0(c,x)=B_0(c,P)=B_0(c,\rho).
 \numberthis\label{eq:rho-field}
\]
Indeed, the first equality follows in both directions from
$P\in\Q(c,x)$ and $x=d^mP^e$, and the second is immediate from
$\rho=s^aP$ with $s\in B_0$.  The coprimality $(a,q)=1$ instead records
that $B_0$ is already generated by the constant used in $\rho$:
if $au+qv=1$, then $s=(s^a)^u d^v$, so $\Q(s^a)=\Q(s)$.

\begin{proposition}[The original-parameter PB cover]
\label{prop:rho-pb}
On the cycle quotient in the original parameter $c$, the map
\[
 \rho:\mathcal C_{d,m}\dashrightarrow\mathbf G_m
\]
has degree
\[
 L_{d,m}=\frac{\nu_d(m)}d
 \numberthis\label{eq:rho-degree}
\]
and satisfies \emph{(PB)} over $B_0$.  Moreover, $c$ separates all
torsion fibers with $x=\rho^e\ne1$.
\end{proposition}

\begin{proof}
The function $P$, and hence $\rho$, has a simple zero at every exact
period-$m$ critical center.  Therefore $\rho$ is not a $p$-th power for
any prime $p$ in the geometric function field.  The binomial criterion,
including its harmless fourth-power clause as in the proof of
Theorem~\ref{thm:unicritical-pb}, proves (PB).

The scaled-parameter degree formula gives
\[
 [B_0(c,x):B_0(x)]=\deg_cM_{d,m}=\frac{e\nu_d(m)}d.
\]
Because $\rho$ is transcendental and $x=\rho^e$, one has
$[B_0(\rho):B_0(x)]=e$.  The tower
\[
 B_0(x)\subset B_0(\rho)\subset B_0(c,\rho)=B_0(c,x)
\]
therefore gives
$[B_0(c,\rho):B_0(\rho)]=\nu_d(m)/d$, proving
\eqref{eq:rho-degree}.  Finally, if the same $c$ occurred in
two torsion fibers, it would give two distinct finite root-of-unity
multiplier cycles for one unicritical polynomial.  The uniqueness
argument in Proposition~\ref{prop:unicritical-separation} excludes
this, and the quotient is \'etale at $x\ne1$.
\end{proof}

Let $H_{d,m}(C,R)$ be a normalized primitive-element equation for
$c$ over $B_0(\rho)$; it has degree $L_{d,m}$ in $C$.  Up to a nonzero
constant independent of $C$,
\[
 M_{d,m}(C,R^e)
 \doteq
 \prod_{\xi\in\mu_e}H_{d,m}(C,\xi R).
 \numberthis\label{eq:original-generic-product}
\]
Indeed, after adjoining $R$ with $R^e=x$, the base change decomposes
into the $e$ components $\rho=\xi R$.  Each is geometrically integral
and has primitive-element equation $H_{d,m}(C,\xi R)$.  These equations
are distinct: an equality
$H_{d,m}(C,R)=H_{d,m}(C,\delta R)$ with
$1\ne\delta\in\mu_e$ would induce a nontrivial automorphism
$c\mapsto c$, $\rho\mapsto\delta\rho$ of
$\overline\Q(c,\rho)=\overline\Q(c,x)$ fixing $c$ and
$x=\rho^e$, which is impossible.  Comparison of degrees in $C$ now
gives \eqref{eq:original-generic-product}.
Define
\[
 \mathcal R_{d,m,k}(C)
 =\Res_X\bigl(\Phi_k(X),M_{d,m}(C,X)\bigr).
 \numberthis\label{eq:original-resultant}
\]
This is, up to the conventional normalization, the original
unicritical delta polynomial, and
\[
 \mathcal R_{d,m,k}(C)
 \doteq \mathcal P_{d,m,k}(d^dC^e).
\]

\begin{theorem}[Original-parameter factorization]
\label{thm:original-parameter-factorization}
Fix $d$ and $m$.  For every sufficiently large $k\ge2$ the following
statements hold.

\begin{enumerate}[label=\textup{(\roman*)}]
\item Over $B_0^c=B_0(\mu_\infty)$, the polynomial
      $\mathcal R_{d,m,k}$ has exactly $e\varphi(k)$ distinct
      irreducible factors, each of degree $L_{d,m}$.
\item Put $K_n=\Q(\zeta_n)$ and $D_n=B_0\cap K_n$.  For every
      \[
       n=kb,\qquad b\mid e,\qquad (k,e/b)=1,
      \]
      the order-$n$ block has $[D_n:\Q]$ irreducible factors over
      $B_0$, each of degree $L_{d,m}[K_n:D_n]$.
\item If $c_0$ lies in the torsion fiber $\rho=\theta$, then
      \[
       B_0(c_0)\cap B_0^c=B_0(\theta).
       \numberthis\label{eq:original-intersection}
      \]
\item If $s\in\Q$---in particular, if $d-1\mid m$---then the complete
      factorization over $\Q$ is indexed by
      \[
       b\mid e,\qquad (k,e/b)=1,
      \]
      and the factor indexed by $b$ has degree
      $L_{d,m}\varphi(kb)$.  Consequently,
      \[
       \mathcal R_{d,m,k}\text{ is irreducible over }\Q
       \quad\Longleftrightarrow\quad
       \operatorname{rad}(d-1)\mid k
       \numberthis\label{eq:original-radical-criterion}
      \]
      for all sufficiently large $k$.
\end{enumerate}
\end{theorem}

\begin{proof}
Apply Theorem~\ref{thm:kummer-torsion-number-field} over $B_0$ to the
PB cover $\rho$, with separating coordinate $c$ and Kummer exponent
$e$.
This proves (i), (iii), and the orbit statement underlying (ii).
Primitive $n$-th roots form $[D_n:\Q]$ orbits under
$\Gal(B_0^c/B_0)$, each of size $[K_n:D_n]$, which proves (ii).  If
$s\in\Q$, the cover and its separating coordinate are defined over
$\Q$; Theorem~\ref{thm:kummer-torsion} and
Corollary~\ref{cor:abstract-factor-count}, with $e$ in place of $r$,
give (iv).
\end{proof}

There is also a complete descent to $\Q$ when $s\notin\Q$, but its
correct indexing is by twisted Kummer orbits rather than by orders
alone.  This distinction prevents a false extension of
\eqref{eq:original-radical-criterion}.

\begin{proposition}[Twisted orbit descent]
\label{prop:twisted-orbit-descent}
Put $\Omega=B_0^c$.  Then $\Omega/\Q$ is Galois.  For
$\sigma\in\Gal(\Omega/\Q)$ set
\[
 \kappa_\sigma=\frac{\sigma(s)}s\in\mu_q.
\]
On
\[
 S_e(k)=\{\theta\in\mu_\infty:\ord(\theta^e)=k\}
\]
define
\[
 \sigma*\theta
 =\kappa_\sigma^{-a}\sigma(\theta).
 \numberthis\label{eq:twisted-action}
\]
For all sufficiently large $k$, the irreducible factors of
$\mathcal R_{d,m,k}$ over $\Q$ are indexed by the orbits
$\mathcal O$ of this action, and the corresponding factor has degree
\[
 L_{d,m}|\mathcal O|.
 \numberthis\label{eq:twisted-degree}
\]
\end{proposition}

\begin{proof}
The field $\Omega$ contains all conjugates $\xi s$ with
$\xi\in\mu_q$, so it is normal over $\Q$.  The quantities
$\kappa_\sigma$ form the usual Kummer cocycle:
\[
 \kappa_{\sigma\tau}
 =\kappa_\sigma\,\sigma(\kappa_\tau).
\]
Thus \eqref{eq:twisted-action} is a group action.  Since $q\mid e$,
\[
 (\sigma*\theta)^e=\sigma(\theta^e),
\]
so the action preserves $S_e(k)$.

The irreducible $\Omega$-factors are the fibers $\rho=\theta$.  If a
point $P_0$ satisfies $\rho(P_0)=\theta$, then
\[
 \rho(\sigma P_0)
 =s^a\sigma(P(P_0))
 =\kappa_\sigma^{-a}\sigma\bigl(s^aP(P_0)\bigr)
 =\kappa_\sigma^{-a}\sigma(\theta).
\]
Thus
$\Gal(\Omega/\Q)$ permutes the distinct irreducible $\Omega$-factors
exactly by \eqref{eq:twisted-action}.  Galois descent says that the
product over one orbit belongs to $\Q[C]$.  It is irreducible: after
extension to $\Omega$, a proper factor over $\Q$ would correspond to a
nonempty proper $\Gal(\Omega/\Q)$-stable subset of that orbit.  Every
$\Q$-factor arises in this way.  Its degree is
\eqref{eq:twisted-degree}.
\end{proof}

\begin{remark}
The twisted action is genuinely necessary.  For $d=3$, $m=1$, and a
primitive cubic multiplier, direct elimination gives
\[
 729C^4+27C^2+169.
\]
Over $\Q(\sqrt3)$ its two torsion-order blocks are conjugate and merge
over $\Q$.  Thus the simple order indexing valid when $s\in\Q$ cannot
be asserted in general.
\end{remark}
